% ============================================================
% Capítulo 6 — Visual Network: conceptos + código
% ============================================================
\chapter{Red visual: EfficientNet-B0 y convolución separable}
\label{cap:visual-network}

Este capítulo combina \textbf{explicaciones conceptuales} de las
piezas usadas en la rama visual con el \textbf{código línea por
línea} del proyecto.

\section{¿Qué es EfficientNet?}

\subsection{El problema}

Hasta 2019, la mayoría de las CNNs se escalaban de forma
\emph{arbitraria}: o se hacían más profundas (más capas), o más
anchas (más canales), o con mayor resolución de entrada. Cada
familia tenía su propia receta y nadie sabía cuál era la mejor.

\subsection{Compound scaling}

\eff{}~\cite{tan2019efficientnet} introdujo una idea elegante:
\textbf{escalar las tres dimensiones simultáneamente} con un
ratio fijo ($\phi$). El paper demostró que esto da mejor
accuracy por FLOP que cualquier escalado aislado.

\begin{equation}
  \begin{aligned}
    \text{depth}   &= \alpha^{\phi} \\
    \text{width}   &= \beta^{\phi} \\
    \text{resolution} &= \gamma^{\phi} \\
    \text{con} \quad & \alpha \cdot \beta^2 \cdot \gamma^2 \approx 2
  \end{aligned}
  \label{eq:compound-scaling}
\end{equation}

\noindent La familia EfficientNet va desde B0 (la más chica) hasta
B7 (la más grande). En este proyecto se usa \textbf{B0} porque es
la única que cabe en hardware modesto con un budget razonable.

\subsection{¿Qué es MBConv?}

MBConv (\emph{Mobile Inverted Bottleneck Convolution}) es el
bloque básico de EfficientNet. Inventado por MobileNetV2 y
mejorado por EfficientNet. Tiene tres etapas:

\begin{enumerate}
  \item \textbf{Expansión}: conv 1x1 que aumenta canales
        ($\times 6$ típicamente, por eso "MBConv6").
  \item \textbf{Depthwise conv}: conv $k \times k$ por canal
        (máscara espacial).
  \item \textbf{Proyección}: conv 1x1 que reduce canales.
\end{enumerate}

\noindent A diferencia del bottleneck clásico de ResNet, MBConv
\emph{expande} primero y \emph{comprime} después ("inverted").

\subsection{Layout de B0}

\begin{table}[H]
\centering
\caption{Las 9 etapas de EfficientNet-B0.}
\label{tab:effnet-stages-concepto}
\begin{tabular}{clrrr}
\toprule
\textbf{Stage} & \textbf{Operador} & \textbf{Resolución} & \textbf{Canales} & \textbf{Layers} \\
\midrule
1  & Conv3x3       & $224 \times 224$ & 32   & 1 \\
2  & MBConv1, k3x3 & $112 \times 112$ & 16   & 1 \\
3  & MBConv6, k3x3 & $112 \times 112$ & 24   & 2 \\
4  & MBConv6, k5x5 & $56 \times 56$   & 40   & 2 \\
5  & MBConv6, k3x3 & $28 \times 28$   & 80   & 3 \\
6  & MBConv6, k5x5 & $14 \times 14$   & 112  & 3 \\
7  & MBConv6, k5x5 & $14 \times 14$   & 192  & 4 \\
8  & MBConv6, k3x3 & $7 \times 7$     & 320  & 1 \\
9  & Conv1x1+Pool  & $7 \times 7$     & 1280 & 1 \\
\bottomrule
\end{tabular}
\end{table}

\noindent\textbf{Cada etapa divide la resolución a la mitad}
(excepto 2 y 5 que mantienen). La \textbf{etapa 9} es la que se usa
en TextReIDNet: produce 1280 canales a $7 \times 7$ (con entrada
$224 \times 224$) o $16 \times 16$ (con entrada $512 \times 512$).

\section{¿Qué es la convolución separable en profundidad (DSC)?}

\subsection{El problema}

Una convolución 2D estándar $k \times k$ sobre una imagen con $C$
canales de entrada y $C'$ canales de salida tiene:
\begin{itemize}
  \item \textbf{Parámetros}: $k \cdot k \cdot C \cdot C'$.
  \item \textbf{FLOPs}: $k^2 \cdot C \cdot C' \cdot H \cdot W$.
\end{itemize}

\noindent Para $k=3$, $C=C'=1280$: \textbf{14,7M} parámetros y
muchos FLOPs.

\subsection{La factorización}

La \emph{depthwise separable convolution} (DSC)~\cite{howard2017mobilenets}
factoriza esa operación en dos:

\begin{enumerate}
  \item \textbf{Depthwise}: una convolución $k \times k$ por canal,
        con $k^2 \cdot C$ parámetros (en vez de $k^2 \cdot C \cdot C'$).
  \item \textbf{Pointwise}: una convolución $1 \times 1$ para mezclar
        canales, con $C \cdot C'$ parámetros.
\end{enumerate}

\noindent\textbf{Resultado}: $k^2 \cdot C + C \cdot C'$ parámetros,
una reducción de $\sim k^2$ veces. Para $k=3$: $\sim 9\times$.

\subsection{Visualización}

\begin{figure}[H]
\centering
\begin{tikzpicture}[
    node distance=0.4cm and 0.6cm,
    every node/.style={font=\sffamily\footnotesize, align=center},
    box/.style={
        rectangle, draw=black, thick, rounded corners=2pt,
        minimum height=1.2cm, minimum width=1.6cm, fill=blue!10
    },
    arr/.style={-{Stealth[length=2.5mm,width=2mm]}, thick}
]
% Standard
\node[box] (in1) {Input\\$C$};
\node[box, right=of in1] (k1) {Conv $k{\times}k$\\$C \to C'$};
\node[box, right=of k1] (out1) {Output\\$C'$};
\draw[arr] (in1) -- (k1); \draw[arr] (k1) -- (out1);
\node[above=0.3cm of k1] {Estándar};

% DSC
\node[box, below=2cm of in1] (in2) {Input\\$C$};
\node[box, right=of in2] (d) {Depthwise\\$k{\times}k$\\por canal};
\node[box, right=of d] (p) {Pointwise\\$1{\times}1$\\$C \to C'$};
\node[box, right=of p, minimum width=1.6cm] (out2) {Output\\$C'$};
\draw[arr] (in2) -- (d); \draw[arr] (d) -- (p); \draw[arr] (p) -- (out2);
\node[above=0.3cm of d] {DSC};

\node[below=1.5cm of p, font=\small, text width=12cm, align=center]
    {Parámetros: estándar $= k^2 \cdot C \cdot C'$ \\[2pt]
    DSC $= k^2 \cdot C + C \cdot C'$ \\[-2pt]
    (reducción $\approx k^2$ veces)};
\end{tikzpicture}
\caption{Comparación entre convolución estándar y DSC.}
\label{fig:dsc}
\end{figure}

\section{¿Qué es AdaptiveMaxPool2d?}

Un \emph{max pooling adaptativo} recibe como argumento el tamaño
\emph{deseado de salida}, no el stride ni el kernel. Calcula
automáticamente los hiperparámetros para que el tensor de salida
tenga exactamente esa forma.

\begin{itemize}
  \item \texttt{nn.AdaptiveMaxPool2d((1, 1))}: salida siempre
        $(B, C, 1, 1)$.
  \item Para cada canal $c$, devuelve el valor máximo sobre todas
        las posiciones espaciales.
  \item Equivale a un \emph{GAP} (global max pooling) si el kernel
        efectivo cubre toda la imagen.
\end{itemize}

\noindent\textbf{¿Por qué?} Después del EfficientNet tenemos un
tensor $(B, 1024, 16, 16)$. Necesitamos un vector por imagen. El
max pooling global colapsa $(16, 16) \to (1, 1)$ quedándose con la
activación más fuerte de cada canal.

\section{Código: la clase \texttt{VisualNetwork}}

\begin{minted}[language=Python, numbers, firstnumber=7, linerange=7-39]{python}
class VisualNetwork(nn.Module):
    """Backbone visual: EfficientNet-B0 preentrenado."""

    def __init__(self):
        # Constructor de nn.Module (registra parámetros)
        super(VisualNetwork, self).__init__()

        # Cargar EfficientNet-B0 con pesos preentrenados en ImageNet
        # EfficientNet_B0_Weights.DEFAULT = mejores pesos disponibles
        # (~5.3M parámetros, input esperado 224x224)
        self.model = models.efficientnet_b0(
            weights=EfficientNet_B0_Weights.DEFAULT
        )

        # `model.features` es un nn.Sequential con 9 bloques
        # (ver Tabla 6.1). Vamos a iterar bloque por bloque.
        self.feature_blocks = self.model.features

    def forward(self, x):
        """x: (B, 3, H, W) imagen preprocesada.
           returns: (B, 1280, H/32, W/32) salida etapa 9."""
        # Lista para acumular features de cada bloque
        features = []

        # Iterar sobre los 9 bloques secuenciales de EfficientNet-B0
        for block in self.feature_blocks:
            # Aplicar bloque (conv + BN + activation)
            x = block(x)

            # Guardar feature de este bloque
            features.append(x)

        # Devolver SOLO la salida del último bloque (etapa 9)
        # = 1280 canales, resolución espacial /32
        return features[-1]
\end{minted}

\section{Código: \texttt{DepthwiseSeparableConv}}

\begin{minted}[language=Python, numbers, firstnumber=259, linerange=259-287]{python}
class DepthwiseSeparableConv(nn.Module):
    """Convolución separable en profundidad: depthwise + pointwise."""

    def __init__(self, in_channels=None, out_channels=None,
                 kernel_size=3, bias=True, is_final_conv=False):
        # Constructor de nn.Module
        super(DepthwiseSeparableConv, self).__init__()

        # Depthwise convolution: una convolución k×k por canal
        # groups=in_channels hace que cada filtro actúe sobre UN canal
        # padding=1 mantiene la dimensión espacial
        self.depthwise_convolution = nn.Conv2d(
            in_channels=in_channels,
            out_channels=in_channels,    # mismo número de canales
            kernel_size=kernel_size,
            groups=in_channels,          # ← clave: depthwise
            bias=bias,
            padding=1                    # mantiene H, W
        )

        # Pointwise convolution: 1×1 que mezcla canales
        # in_channels -> out_channels
        self.pointwise_convolution = nn.Conv2d(
            in_channels=in_channels,
            out_channels=out_channels,   # ← dimensión de salida
            kernel_size=1,
            groups=1,                    # convolución normal
            bias=bias
        )

        # Encadenar las dos operaciones
        self.depthwise_separable_conv = nn.Sequential(
            self.depthwise_convolution,
            self.pointwise_convolution
        )

        # RetinaNet-style bias init (no se usa en TextReIDNet)
        self.is_final_conv = is_final_conv
        if bias and is_final_conv:
            self.is_final_conv = True

        # Inicialización Kaiming (para Conv2d y Linear)
        self.depthwise_separable_conv.apply(weights_init_kaiming)

    def forward(self, x):
        return self.depthwise_separable_conv(x)
\end{minted}

\section{Etapas con shapes exactas (input 512×512)}

\begin{table}[H]
\centering
\caption{Etapas con shapes para input $512 \times 512$ en este proyecto.}
\label{tab:effnet-shapes}
\begin{tabular}{clrrr}
\toprule
\textbf{Stage} & \textbf{Operador} & \textbf{Output HxW} & \textbf{Output Ch} & \textbf{Layers} \\
\midrule
0  & Conv3x3       & $256 \times 256$  & 32   & 1 \\
1  & MBConv1, k3x3 & $256 \times 256$  & 16   & 1 \\
2  & MBConv6, k3x3 & $128 \times 128$  & 24   & 2 \\
3  & MBConv6, k5x5 & $64 \times 64$    & 40   & 2 \\
4  & MBConv6, k3x3 & $32 \times 32$    & 80   & 3 \\
5  & MBConv6, k5x5 & $16 \times 16$    & 112  & 3 \\
6  & MBConv6, k5x5 & $16 \times 16$    & 192  & 4 \\
7  & MBConv6, k3x3 & $8 \times 8$      & 320  & 1 \\
8  & Conv1x1       & $8 \times 8$      & 1280 & 1 \\
\bottomrule
\end{tabular}
\end{table}

\section{Flujo de tensores}

\begin{minted}[language=Python, basicstyle=\ttfamily\small, frame=single]{python}
# Entrada: imagen preprocesada (normalizada)
x:  (B, 3, 512, 512)

# EfficientNet-B0
x = visual_network(x)
# -> (B, 1280, 16, 16)  <- etapa 9, depende del tamaño de entrada

# DSC 1280 -> 1024
x = visual_features_downscale(x)
# -> (B, 1024, 16, 16)

# AdaptiveMaxPool2d((1, 1))
x = adaptive_max_pooling(x)
# -> (B, 1024, 1, 1)

# DSC final 1024 -> 1024
x = depthwise_seperable_convolution(x)
# -> (B, 1024, 1, 1)

# squeeze última dim
x = x.squeeze(-1).contiguous()
# -> (B, 1024, 1)
\end{minted}

\section{Parámetros}

\begin{table}[H]
\centering
\caption{Parámetros por componente visual.}
\label{tab:vis-params-linea}
\begin{tabular}{lr}
\toprule
\textbf{Componente} & \textbf{Parámetros} \\
\midrule
\texttt{EfficientNet-B0 features}   & \textasciitilde 5,29 M \\
\texttt{DSC(1280, 1024)}             & \textasciitilde 14 K \\
\texttt{AdaptiveMaxPool2d(1,1)}      & 0 \\
\texttt{DSC(1024, 1024)}             & \textasciitilde 10 K \\
\midrule
\textbf{Total visual}                & \textbf{\textasciitilde 5,31 M} \\
\bottomrule
\end{tabular}
\end{table}
